\section*{Exercises}
\addcontentsline{toc}{section}{Exercises}

\subsection*{A. Theory}

\begin{exercise}
\exerciseid{9.A1}
Explain why solving normal equations can be less stable than solving least
squares with QR. Your answer should mention $A^TA$.
\end{exercise}

\begin{exercise}
\exerciseid{9.A2}
Suppose $Q$ has orthonormal columns. Prove that the least-squares solution to
$Qx=b$ is $\hat{x}=Q^Tb$.
\end{exercise}

\subsection*{B. By Hand}

\begin{exercise}
\exerciseid{9.B1}
Fit a line $y=a+bt$ through the points $(0,1)$, $(1,2)$, and $(2,2)$ in the
least-squares sense.
\end{exercise}

\begin{exercise}
\exerciseid{9.B2}
Apply one full Gram--Schmidt computation to
\[
  v_1=(1,1,0),\qquad v_2=(1,0,1).
\]
Find an orthonormal basis for $\Span(v_1,v_2)$.
\end{exercise}

\subsection*{C. Python / NumPy}

\begin{exercise}
\exerciseid{9.C1}
Build a Vandermonde matrix for five sample points and compare the
normal-equation solution, a QR-based solution, and
\texttt{np.linalg.lstsq}.
\end{exercise}

\begin{exercise}
\exerciseid{9.C2}
Implement classical Gram--Schmidt using only basic \NumPy operations. Test
whether your output satisfies \texttt{Q.T @ Q = I} and \texttt{Q @ R = A}.
\end{exercise}

\subsection*{D. Challenge}

\begin{exercise}
\exerciseid{9.D1}
Explain why increasing polynomial degree can reduce training error while
increasing validation error. Use the column-space picture in your explanation.
\end{exercise}
